\section{A word rule with no vertical period}
\label{sec:word}

We begin with a finite word constraint whose solutions on an infinite strip
cannot have a vertical period if every column is nonconstant.  The constraint
is the two-digit $p$-adic Sudoku rule of
\citet[Definition~4.5]{GTUndecidable}.  The affine-square and rescaling
arguments below follow the framework of
\citet[Theorem~4.7, Lemma~4.8, and Section~4.2]{GTUndecidable}; we give the
complete argument for the particular obstruction needed here.

Fix a prime $p>200$, and set
\[
    N=p^2,\qquad \Sigma=\Fp^\times,
    \qquad I=\{0,1,\ldots,N-1\}.
\]
When an integer occurs in an expression over $\Fp$, its residue modulo $p$
is understood.

\begin{definition}[Allowed words and the line rule]
\label{word:definition}
A word $w\colon I\to\Sigma$ is \emph{allowed} if there are integers $a,b$,
not both divisible by $p$, such that, for every $n\in I$,
\[
 w(n)=
 \begin{cases}
   an+b\pmod p, & p\nmid an+b,\\[2pt]
   (an+b)/p\pmod p, & p\mid an+b\text{ and }p^2\nmid an+b.
 \end{cases}
\]
At indices with $p^2\mid an+b$, the value may be any member of $\Sigma$.
An array $W\colon I\times\Z\to\Sigma$ satisfies the \emph{line rule} if
\[
    n\longmapsto W(n,dn+e)
\]
is an allowed word for every $d,e\in\Z$.
A \emph{positive vertical period} of $W$ is an integer $M>0$ such that
$W(n,m+M)=W(n,m)$ for all $(n,m)\in I\times\Z$.
\end{definition}

The allowed-word constraint depends only on the residues of $a,b$ modulo
$p^2$.  Indeed, these residues determine the divisibility alternatives
and each of the prescribed values.  In particular, it is a finite
constraint on words over the finite alphabet $\Sigma$.

For every allowed word, reduction of a witnessing pair gives a nonzero
affine form $\ell(n)=\bar a n+\bar b$ over $\Fp$.  The word agrees with
$\ell$ wherever $\ell\ne0$.  A nonzero affine form has at most one zero
residue class, so there are at most $N/p=p$ exceptional indices on $I$.
Here and below, ``nonzero affine form'' means that its coefficients are
not all zero; a nonzero constant form is permitted.

\begin{lemma}[Global affine approximation]
\label{word:affine}
If $W$ satisfies the line rule, there are $A,B,C\in\Fp$, not all zero,
such that
\[
    W(n,m)=An+Bm+C
    \quad\text{whenever }An+Bm+C\ne0.
\]
\end{lemma}

\begin{proof}
For each line $m=dn+e$ with $d\in\{-1,0,1\}$, choose the nonzero affine
form $\ell_{d,e}(n)$ supplied by an allowed-word witness.  Call an index on
that line exceptional when $\ell_{d,e}(n)=0$.

First we find a consecutive $4\times4$ square containing no exceptional
cell for any of these three line directions.  In the rectangle
$I\times\{0,\ldots,K-1\}$ there are $(N-3)(K-3)$ such squares.  The
numbers of lines of slopes $0,1,-1$ meeting the rectangle are respectively
$K,K+N-1,K+N-1$.  Each line has at most $p$ exceptional cells on the
whole strip, and each cell belongs to at most $16$ candidate squares.
Consequently at most
\[
    16p(3K+2N-2)\leq16p(3K+2N)
\]
squares are excluded.  Since
\[
    N-3-48p=p^2-48p-3>0,
\]
we can choose an integer $K\geq4$ with
\[
    K>\frac{3(N-3)+32pN}{N-3-48p}.
\]
For this choice,
\[
    (N-3)(K-3)-16p(3K+2N)
      =K(N-3-48p)-3(N-3)-32pN>0.
\]
A square of the required kind therefore exists.

Write its column indices as $r,r+1,r+2,r+3$ and its row indices as
$q,q+1,q+2,q+3$.  Each row agrees on the square with an affine function
$U_m n+V_m$.  For $m=q+1,q+2$ and $n=r+1,r+2$, the second difference
along either diagonal vanishes.  Thus, for each sign,
\[
 \begin{split}
  0={}&n(U_{m+1}-2U_m+U_{m-1})
       +(V_{m+1}-2V_m+V_{m-1})\\
     &\qquad\qquad{}\pm(U_{m+1}-U_{m-1}).
 \end{split}
\]
Subtracting the identities at the two consecutive values of $n$ gives
$U_{m+1}-2U_m+U_{m-1}=0$.  Subtracting the two signs gives
$U_{m+1}=U_{m-1}$, because $2\ne0$ in $\Fp$.  Together these imply
$U_{m-1}=U_m=U_{m+1}$.  The two overlapping triples show that all four
row slopes have a common value $A$.  The remaining identities say that
the four intercepts have zero second differences, so they have the form
$V_m=Bm+C$.  Therefore the form
\[
    H(n,m)=An+Bm+C
\]
agrees with $W$ on the square.  It is not the zero form, because every
entry of $W$ lies in $\Sigma$.

Call a cell \emph{good} if either $H(n,m)=0$ or $W(n,m)=H(n,m)$.  We claim
that four consecutive good cells on a line of slope $0,1$, or $-1$ make
the entire line good.  The restriction of $H$ to that line is an affine
form in $n$.  If it is identically zero, the claim follows from the
definition of goodness.  Otherwise, it and the chosen form
$\ell_{d,e}$ each have at most one zero among the four indices, which
are distinct modulo $p$.  At least two of the four indices remain where
both forms are nonzero.  At these indices, goodness and the word rule
identify both forms with $W$.  Two distinct residues determine an
affine form, so the two forms are equal.  Wherever this common form is
nonzero the word rule gives $W=H$, and its zeros are good by definition.
This proves the claim.

Applying the claim to the four square rows makes those entire rows good.
To adjoin a row immediately above a block of four good rows, consider
its cell in column $n$.  At least one of
\[
    n+4\leq N-1\qquad\text{or}\qquad n-4\geq0
\]
holds, since $N\geq8$.  Choose $\varepsilon\in\{-1,1\}$ so that the four
columns $n+\varepsilon,n+2\varepsilon,n+3\varepsilon,n+4\varepsilon$
are in $I$.  If the new row has index $m$, the cells
\[
    (n+j\varepsilon,m-j),\qquad 1\leq j\leq4,
\]
are four consecutive good cells on a diagonal through $(n,m)$.  The
claim makes $(n,m)$ good.  For a row immediately below a block of four
good rows, use $(n+j\varepsilon,m+j)$ instead.  These formulas include
both edge columns of the strip.  Successively adjoining rows in both
directions makes every cell good, proving the lemma.
\end{proof}

\begin{lemma}[Vertical nondegeneracy]
\label{word:vertical}
Let $H(n,m)=An+Bm+C$ be any affine form supplied by
\cref{word:affine}.  Every column of $W$ is nonconstant if and only if
$B\ne0$.  In that case every positive vertical period of $W$ is
divisible by $p$.
\end{lemma}

\begin{proof}
If $B=0$, then $An+C$ is not identically zero.  Some $n\in\{0,\ldots,p-1\}$
therefore has $An+C\ne0$, and the entire column at that $n$ has the
constant value $An+C$.  Conversely, if $B\ne0$, then for any fixed $n$
the form $An+Bm+C$ runs through $\Fp$ as $m$ runs through a complete
residue system modulo $p$.  At its $p-1$ nonzero values, the entries of
$W$ equal those values.  Thus every column takes every symbol of
$\Sigma$, and in particular is nonconstant.

Suppose now that $M>0$ satisfies $W(n,m+M)=W(n,m)$ for every $(n,m)$.
For a fixed $n$, choose $m$ so that both $H(n,m)$ and $H(n,m+M)$ are
nonzero.  There are at most two excluded residue classes, so this is
possible.  Equality of the two entries gives
\[
    BM=H(n,m+M)-H(n,m)=0\quad\text{in }\Fp.
\]
Since $B\ne0$, this implies $p\mid M$.
\end{proof}

\begin{proposition}[The vertical-period obstruction]
\label{word:obstruction}
There is no array $W\colon I\times\Z\to\Sigma$ satisfying the line rule,
having every column nonconstant, and having a positive vertical period.
\end{proposition}

\begin{proof}
Suppose otherwise, and choose the smallest positive integer $M$ that
occurs as a vertical period of any such array.  Choose an array with
period $M$ and an affine approximation $An+Bm+C$.  By
\cref{word:vertical}, $B\ne0$ and $p\mid M$.

Choose integers $u,v$ with $A+Bu=C+Bv=0$ in $\Fp$, and replace $W$ by
\[
    \widetilde W(n,m)=W(n,m+un+v).
\]
On a line of slope $d$ and intercept $e$, this samples the original
line of slope $d+u$ and intercept $e+v$, so the line rule is preserved.
Each column is simply translated in its row coordinate; hence column
nonconstancy and the set of vertical periods are preserved.  The new
affine approximation is $Bm$.  Rename this normalized array $W$.

Define
\[
    W_1(n,m)=W(n,pm).
\]
This array also satisfies the line rule: its line with slope $d$ and
intercept $e$ is the old line with slope $pd$ and intercept $pe$.
By \cref{word:affine}, it has a nonzero affine approximation
$H_1(n,m)=A_1n+B_1m+C_1$.  We next prove $B_1\ne0$, which is the step
needed to retain column nonconstancy under this rescaling.

If $B_1=0$, choose $0\leq n_0<p$ with
$c=A_1n_0+C_1\ne0$.  Apply the word rule for $W$ to the line
$m=n-n_0$, and let $a,b$ be a witnessing pair.  Away from the at most
two zero residue classes of $\bar a n+\bar b$ and $B(n-n_0)$, both
forms equal the array entry on that line.  Since $p-2\geq2$, they agree
at two distinct residues and therefore are the same affine form.  In
particular,
\[
    a\equiv B\pmod p,\qquad b\equiv-Bn_0\pmod p.
\]
It follows that $p\mid an_0+b$ and $p\nmid a$.

For $0\leq k<p$, put $n=n_0+pk$; all these indices lie in $I$,
including the largest possible value $p^2-1$.  At the corresponding
point of the line, the entry is
\[
    W(n_0+pk,pk)=W_1(n_0+pk,k)=c,
\]
where the last equality follows from the nonzero value
$H_1(n_0+pk,k)=c$.  On the other hand,
\[
    a(n_0+pk)+b
       =p\left(\frac{an_0+b}{p}+ak\right).
\]
Whenever the parenthesized expression is nonzero modulo $p$, the second
case of the allowed-word rule forces this entry to equal
\[
    \frac{an_0+b}{p}+ak\pmod p.
\]
As $k$ runs from $0$ to $p-1$, these residues run through all of $\Fp$,
because $p\nmid a$.  Thus $p-1$ of the entries are forced to take the
$p-1$ distinct nonzero values, contradicting their common value $c$.
The remaining index, where the rule prescribes no value, is not needed
for this contradiction.

We have proved $B_1\ne0$, so every column of $W_1$ is nonconstant by
\cref{word:vertical}.  But $M/p$ is a positive vertical period of $W_1$,
since
\[
    W_1(n,m+M/p)=W(n,pm+M)=W(n,pm)=W_1(n,m).
\]
This contradicts the choice of $M$.
\end{proof}

\begin{lemma}[A model for the line rule]
\label{word:model}
For $t\in\Z\setminus\{0\}$, let $\operatorname{ord}_p(t)$ be the
largest nonnegative integer $r$ for which $p^r\mid t$, and define
\[
    f(t)=\frac{t}{p^{\operatorname{ord}_p(t)}}\pmod p,
    \qquad f(0)=1.
\]
Then $f\colon\Z\to\Sigma$, and the array $W(n,m)=f(m)$ satisfies the
line rule and has every column nonconstant.
\end{lemma}

\begin{proof}
The definition gives a nonzero residue for every nonzero integer,
including negative integers.  For every $r\geq0$ and $t\in\Z$ we have
$f(p^rt)=f(t)$; this also holds at $t=0$ by definition.

Fix integers $d,e$.  If $(d,e)\ne(0,0)$, let $r$ be the largest integer
such that $p^r$ divides both coefficients, and write $d=p^ra$, $e=p^rb$.
The pair $a,b$ is not both divisible by $p$, and
\[
    f(dn+e)=f(an+b)\qquad(n\in I).
\]
When $an+b$ has $p$-adic order zero or one, this is exactly the
respective value prescribed by the allowed-word rule.  When
$p^2\mid an+b$, including $an+b=0$, the rule leaves the value
unrestricted.  Thus $a,b$ witness that the word is allowed.  If
$(d,e)=(0,0)$, the word is constantly $1$ and has witness $a=0,b=1$.
Finally $f(1)=1$ and $f(2)=2$, so every column is nonconstant.
\end{proof}

For the encoding that follows, parametrize the lines by
$x=(x_1,x_2)\in\Z^2$ and write
\begin{equation}
\label{word:line-maps}
    L_n(x)=x_2+nx_1,\qquad u_n=(1,-n)\qquad(n\in I).
\end{equation}
The map $L_n\colon\Z^2\to\Z$ is surjective and its kernel is generated
by $u_n$: the equation $x_2+nx_1=0$ is equivalent to
$x=x_1(1,-n)$.  The model supplies the allowed word
$\bigl(f(L_n(x))\bigr)_{n\in I}$ at every $x$.  Moreover, if
$x_1\equiv0$ and $x_2\equiv t\pmod p$ for $t\in\{1,2\}$, every entry
of this word equals $t$.  Indeed, every $L_n(x)$ then has the nonzero
residue $t$, so its last nonzero digit is $t$.
